\chapter{Correctness and Validation}
\label{ch:correctness-validation}

Correctness means that the generated C implementation agrees with the mathematical meaning of the DSL program for the selected field and configuration. The compiler does not provide a formal proof; it combines reference semantics, randomized tests, and handwritten cross-checks.

\section{Reference Semantics}
\label{sec:reference-semantics}

The reference semantics interprets a DSL program after API-level decoding turns byte buffers into field-element sequences. Index expressions denote integers. Buffer expressions such as \(message[i]\) denote the \(i\)-th decoded field element. Field operations denote addition, multiplication, and squaring in \(\mathbb{F}_{2^\pi-\theta}\).

Reductions have their usual mathematical meanings: finite summation, Horner recurrence, or left fold over an explicitly named accumulator. A hash function denotes the final field element exported to the output buffer.

This semantics is independent of limb representations, carry schedules, and vectorization. Generated code is correct if it produces the same field element.

\section{Automatic Boost Tests}
\label{sec:automatic-boost-tests}

For automatic validation, the compiler emits a C++ test file using Boost Multiprecision~\cite{boost}. The scalar reference evaluates the same DSL expressions using arbitrary-precision integers and reduces modulo the selected prime.

Each test generates randomized key and message data, calls the generated C implementation, calls the Boost reference, and compares exported outputs.

This is not a proof, but it catches errors in lowering, load/store code, carry placement, multiplication routines, full reduction, and byte handling.

\section{Manual Tests}
\label{sec:manual-tests}

Some examples use additional tests.

\subsection{Poly1305}
\label{subsec:poly1305-tests}

The Poly1305 DSL module implements the polynomial part of Poly1305. The complete authenticator adds a second key component \(s\); here the test fixes or omits it and compares the generated polynomial output against OpenSSL~\cite{openssl}.

This validates both field arithmetic and byte-level interpretation of Poly1305 blocks.

\subsection{HKM}
\label{subsec:hkm-tests}

HKM is tested in two forms: a corrected recursive recurrence from the new-designs paper, and an iterative left-fold implementation. The manual test compares both on randomized inputs.

The iterative version is intended for benchmarking because it avoids recursive generated C.

\section{Cross-Checks}
\label{sec:correctness-cross-checks}

Generated code is validated by Boost references, OpenSSL checks, and targeted load/store and arithmetic tests.

The tests run for the selected field, representation, unrolling factor, and delayed-realignment mode.
